\chapter{High Level Requirements}

\section{System}
\begin{enumerate}[leftmargin=\dimexpr\linewidth-13cm-\rightmargin\relax,label={\bfseries{SYS-REQ-\arabic*}}]
    \item Space Packets (except idle space packets) that are send to the System on the encoder input shall be outputted on the decoder output, these packets shall be identical
    \item Space Packets shall be delivered on Ground with the same bit order as in Space provided
    \item The Pseudo Random Number generator shall use the generator polynomial \(h(x) = x^{17} + x^{14} + 1\) (ref. to CCSDS 131.0-B-5 10.4.1)
    \item The Pseudo Random Number generator shall use the starting Condition 1000111000111000 (ref. to CCSDS 131.0-B-5 10.4.3)
    \item The Reed Solomon (R/S) shall use \(F(x) = x^8 + x^7 + x^2 + x + 1\) over a galois field of GF(2) (ref. to CCSDS 131.0-B-5 4.3.3)
    \item The R/S Encoder shall have a error correction capability of 16 symbols per R/S block ($E=16$)
    \item The Convolutional code shall use $G1 = 1111001$ and $G2 = 1011011$ as Connection vectors
    \item The Convolutional code shall use a $50\%$ Code rate
    \item There shall be NO interleaving used in the Reed Solomon generator
\end{enumerate}

\section{Space}
\begin{enumerate}[leftmargin=\dimexpr\linewidth-13cm-\rightmargin\relax,label={\bfseries{SPC-REQ-\arabic*}}]
    \item The system shall not loose space packets, except idle space packets, in the processing chain
    \item The system shall be able to produce at least 2 Mbit/second of telemetry (TM) data, on the output of the 131 encoder
    \item The system shall interface with the Processing System (PS) with an advanced extensible interface (AXI)
    \item The AXI interface between PS and system shall be clocked to fulfill the data throughput requirements
    \item The system shall interface with the antenna system using an AXI interface
    \item The system shall be implemented on a ZYNQ 7000 (PYNQ Z2) system on a chip (SoC)

\subsection{CCSDS-132.0-B-3}
    \item The TM transfer frame length shall be exactly 2046 octets (\href{https://nc.seesat.eu/f/205229}{SeeSat-1-CCSDS-Stack})
    \item The system shall be CCSDS-132.0-B-3 compliant
    \item The operational control field (OCF) should be implemented
    \item The OCF should be configurable on synthesis
    \item The frame error control field (FECF) should be implemented
    \item The FECF should be configurable on synthesis
    \item The system shall be able to pack partial space packets into TM transfer frames
    \item The system shall be able to calculate the payload trailer of the TM transfer frames without the space data link security (SDLS) attachment
    \item The system shall be able to send data to the TM SYNCHRONIZATION AND CHANNEL CODING layer (CCSDS 131.0-B-5)
    \item The system shall be able to receive data from multiple sources (Virtual Channels)
    \item The system should implement multiple virtual channels
    \item The system should implement multiple master channels
    \item The system should implement a secondary frame header
    \item The secondary frame header should be configurable

\subsection{CCSDS-131.0-B-5}
    \item The system shall use concatenated coding scheme
    \item The system shall use Reed Solomon coding for outer  codes
    \item The system shall use pseudo randomization
    \item The system shall attach a 32-Bit attached synchronization markers (ASM) to the data stream
    \item The system shall use convolutional coding for inner coding
    \item The system shall expose its configuration, in accordance with configuration management
\end{enumerate}


\section{Ground}
\begin{enumerate}[leftmargin=\dimexpr\linewidth-13cm-\rightmargin\relax,label={\bfseries{GND-REQ-\arabic*}}]
    \item The system shall be able to indicate to the user when a packet cannot be corrected
    \item The system shall be able to correct errors that are correctable by the algorithms used
    \item The system shall interface with further systems via TCP packets
    \item Transmission of correct(ed) packets to further Mission Control systems shall only include the binary data of the space packet
    \item The ground system shall be implemented on a ZYNQ 7000 (PYNQ Z2) SoC

\subsection{CCSDS\_131.0-B-5}
    \item The system shall be connected to the receiving antenna via Ethernet/TCP or AXI Stream.
    \item The system shall receive TM data from the antenna as raw bits
    \item The system shall decode these raw bits using an approprate decoding algorithm for convolutional coding    
    \item The system shall subtract the injected pseudo random values from the data returned by the asm decoder
    \item The ground system shall verify the received bit count using the transmitted Attached Synchronization Marker (ASM)    
    \item The system shall decode the Reed Solomon codes
    \item The data generated by this layer (CCSDS 131) shall be transmitted to CCSDS 132


\subsection{CCSDS\_132.0-B-3}
    \item The system shall extract the space data packets
    \item The extracted space data packets shall be send via TCP to the next systems
    \item The system shall de-multiplex the space packet to its corresponding Virtual- and Master channel
    \item The system should implement the decoding of the FECF
    \item The FECF should be configurable on synthesis
    \item The system should implement the decoding of the OCF
    \item The OCF should be configurable on synthesis
    \item The system should implement a secondary frame header
    \item The secondary frame header should be configurable
    \item If data is not correctable by the CCSDS 131, uncorrected data shall still be received by CCSDS 132 and sent out to higher protocol layers if possible
\end{enumerate}


\section{Non Functional}
\begin{enumerate}[leftmargin=\dimexpr\linewidth-13cm-\rightmargin\relax,label={\bfseries{NFUNC-REQ-\arabic*}}]
    \item The systems program code and hardware configuration shall follow a clear and defined standard (See VHDL Standard - SeeSat-TMTC-101)
    \item The systems program code and hardware configuration shall be written in an extendable and maintainable way
\end{enumerate}

